% ============================================================
%  TechSource - MATLAB MCP Server Setup Guide
%  Combined: TOML-config bridges [Method A] + JSON-config / any
%  other bridge [Method B] -- no specific tool is named
% ============================================================
\documentclass[11pt, letterpaper]{article}

\usepackage[top=0.7in, bottom=1.0in, left=0.85in, right=0.85in, footskip=0.35in]{geometry}
\usepackage{graphicx}
\usepackage{xcolor}
\usepackage[T1]{fontenc}
\usepackage{textcomp}
\usepackage{helvet}
\usepackage{microtype}
\usepackage[nobottomtitles*]{titlesec}
\usepackage{titling}
\usepackage{enumitem}
\usepackage{hyperref}
\usepackage{fancyhdr}
\usepackage{lastpage}
\usepackage{booktabs}
\usepackage{array}
\usepackage{parskip}
\usepackage{tcolorbox}
\usepackage{caption}
\usepackage{float}
\usepackage{amssymb}
\usepackage{fancyvrb}

\renewcommand{\familydefault}{\sfdefault}

\definecolor{TSorange}{RGB}{230, 126, 14}
\definecolor{TSgray}{RGB}{100, 100, 100}
\definecolor{TSmaroon}{RGB}{128, 0, 32}
\definecolor{TSdark}{RGB}{30, 30, 30}
\definecolor{TSlightgray}{RGB}{245, 245, 245}
\definecolor{TSblue}{RGB}{30, 90, 160}
\definecolor{TSboxbg}{RGB}{255, 248, 235}
\definecolor{TScodebg}{RGB}{243, 245, 248}
\definecolor{TSframe}{RGB}{185, 185, 185}

\hypersetup{
  colorlinks = true,
  linkcolor  = TSorange,
  urlcolor   = TSblue,
  citecolor  = TSgray,
  pdftitle   = {MATLAB MCP Server Setup},
  pdfauthor  = {TechSource},
}

\titleformat{\section}
  {\color{TSorange}\large\bfseries\sffamily}
  {Step \thesection.}{0.6em}{}
  [{\color{TSorange}\titlerule[0.8pt]}]

\titleformat{\subsection}
  {\normalsize\bfseries\sffamily}
  {}{0em}{}

\titlespacing{\section}{0pt}{10pt}{5pt}
\titlespacing{\subsection}{0pt}{9pt}{4pt}

\renewcommand{\bottomtitlespace}{0.10\textheight}

\pagestyle{fancy}
\fancyhf{}
\setlength{\headheight}{46pt}
\fancyhead[L]{\includegraphics[height=0.42in]{techsource_logo.png}}
\fancyhead[R]{\includegraphics[height=0.42in]{mw-reseller-badge-rgb-def.png}}
\fancyfoot[C]{\color{TSgray}\small Page \thepage\ of \pageref{LastPage}}
\renewcommand{\headrulewidth}{0.4pt}

\fancypagestyle{plain}{\fancyhf{}\renewcommand{\headrulewidth}{0pt}\renewcommand{\footrulewidth}{0pt}}

\captionsetup{
  font={small,color=TSgray},
  labelfont={bf,color=TSorange},
  format=plain,
  justification=centering,
  skip=4pt,
}

\tcbuselibrary{skins, breakable}

\newtcolorbox{notebox}{
  colback=TSboxbg, colframe=TSorange,
  fonttitle=\bfseries\sffamily\small,
  title={$\bigstar$\; Note},
  left=6pt, right=6pt, top=4pt, bottom=4pt, breakable,
  fontupper=\raggedright,
}

\newtcolorbox{assumptionbox}[1]{
  colback=TSboxbg, colframe=TSorange,
  fonttitle=\bfseries\sffamily\small,
  title={#1},
  left=6pt, right=6pt, top=4pt, bottom=4pt, breakable,
  fontupper=\raggedright,
}

\newtcolorbox{codebox}{
  colback=TScodebg, colframe=TSframe,
  boxrule=0.5pt, arc=2pt,
  left=6pt, right=6pt, top=5pt, bottom=5pt, breakable,
  fontupper=\ttfamily\small\raggedright,
}

\newtcolorbox{codeblock}{
  colback=TScodebg, colframe=TSframe,
  boxrule=0.5pt, arc=2pt,
  left=6pt, right=6pt, top=5pt, bottom=5pt, breakable,
}
% Same styling as codeblock, but never splits across a page --- used where the
% whole block needs to stay together (e.g. a script the reader will copy in one go).
\newtcolorbox{codeblocknobreak}{
  colback=TScodebg, colframe=TSframe,
  boxrule=0.5pt, arc=2pt,
  left=6pt, right=6pt, top=5pt, bottom=5pt,
}
\fvset{fontsize=\footnotesize, xleftmargin=0pt, formatcom=\color{TSdark}}

% Method A / Method B labelled dividers -- visually distinct so each step's
% two methods are easy to tell apart at a glance. Neither name names a
% specific product -- the split is by config file shape only.
\newcommand{\methodA}{\subsection*{\color{TSblue}$\blacktriangleright$~Method A --- If the Bridge Uses a TOML Config File}}
\newcommand{\methodB}{\subsection*{\color{TSmaroon}$\blacktriangleright$~Method B --- If It Uses JSON, or Anything Else}}

\setlist[enumerate,1]{label=\textbf{\arabic*.}, leftmargin=*, itemsep=4pt, topsep=4pt}
\setlist[itemize,1]{leftmargin=*, itemsep=3pt, topsep=3pt, label=\textcolor{TSorange}{\textbullet}}

\setlength{\parskip}{5pt}
\setlength{\parindent}{0pt}

\newcommand{\tick}{\textcolor{TSorange}{$\square$}}
\newcommand{\arrow}{\,$\rightarrow$\,}

% ============================================================
\begin{document}
% ============================================================

\begin{center}
  \color{TSorange}\rule{\textwidth}{2pt}\\[5pt]
  {\fontsize{18}{22}\selectfont\bfseries\color{TSdark}
    MATLAB MCP Server Setup Guide}\\[4pt]
  {\small\color{TSgray}
    For Bridges Using a TOML or JSON Configuration File}\\[2pt]
  {\small\color{TSgray}
    Revised 19 August 2026}\\[5pt]
  \color{TSorange}\rule{\textwidth}{2pt}
\end{center}

\vspace{6pt}

% ============================================================
\section*{Overview}
% ============================================================

This guide does not assume which local model bridge is installed on the target machine, and does
not name any specific product. Instead it splits by the \textbf{shape of that bridge's own
configuration file}, since that alone is usually enough to tell which method applies:

\begin{itemize}
  \item {\color{TSblue}\textbf{Method A}} --- for bridges that store configuration as
        \textbf{TOML} (recognizable by a \path{config.toml}-style file with section headers
        written like \path{[[mcp_servers]]}).
  \item {\color{TSmaroon}\textbf{Method B}} --- for bridges that store configuration as
        \textbf{JSON} (recognizable by curly braces and a key such as \path{"mcpServers": \{...\}}),
        or for any bridge whose format is unknown.
\end{itemize}

Open whatever configuration file the bridge already has to see which shape it is, then follow the
matching method within each step below. \textbf{Try Method A first if unsure; if the config file
does not look like TOML, skip straight to Method B.}

\begin{notebox}
  Steps that do not depend on the bridge's configuration format (downloading the files, raising
  the context limit) are given only once, since the two methods would be identical anyway.
\end{notebox}

\textbf{What you will need:}
\begin{itemize}
  \item A Windows PC
  \item A stable internet connection, for the downloads in Step 1 only
  \item Enough free disk space for the MCP server binary and the skill catalogs (well under 1 GB)
\end{itemize}

\textbf{Two assumptions this guide relies on:}

\begin{assumptionbox}{Assumption 1 --- A Bridge Already Exists}
  The target machine already has a working local model connector bridge connecting a local LLM
  to MCP servers. Tool-calling infrastructure is a given --- there is no need to test whether the
  model itself has file or tool access before registering anything.
\end{assumptionbox}

\begin{assumptionbox}{Assumption 2 --- Context Window Is Not a Hard Constraint}
  The context limit of the local model can be increased as needed, so the full MATLAB/Simulink
  tool set and skill content can be loaded in full, without pruning or selective loading.
\end{assumptionbox}


% ============================================================
\section{Download Everything Into One Folder}
% ============================================================

This step is identical either way --- it does not depend on the bridge's configuration format.
Everything is downloaded by hand from the browser and extracted into one working folder. By the
end of this step, that folder looks like this:

\begin{codeblock}
\begin{Verbatim}
C:\mcp-deploy\matlab\
  matlab-mcp-server\          <- MCP server binary            (Step 2 below)
  matlab-agentic-toolkit\     <- skills catalog + templates    (Step 3 below)
  simulink-agentic-toolkit\   <- Simulink tools/skills, optional (Step 4 below)
\end{Verbatim}
\end{codeblock}

\begin{enumerate}
  \item Create the working folder: \path{C:\mcp-deploy\matlab}.

  \item Inside it, create a subfolder named \path{matlab-mcp-server}. Open\\
        \url{https://github.com/matlab/matlab-mcp-server/releases} in a browser, open the
        \textbf{latest} release, download the \texttt{.exe} listed there for your platform, and
        save it into that subfolder.

  \item Go back to \path{C:\mcp-deploy\matlab}, then:
  \begin{enumerate}[label=(\alph*)]
    \item Open \url{https://github.com/matlab/matlab-agentic-toolkit}.
    \item Click the green \textbf{Code} button, then choose \textbf{Download ZIP}.
    \item Right-click the downloaded ZIP and choose \textbf{Extract All}.
    \item Click \textbf{Browse}, set the destination to \path{C:\mcp-deploy\matlab} itself ---
          do not accept whatever folder is pre-filled --- then click \textbf{Extract}.
    \item \textbf{Rename} the extracted \path{matlab-agentic-toolkit-main} folder to
          \path{matlab-agentic-toolkit}.
  \end{enumerate}

  \item Only if Simulink tools/skills are also needed, still inside \path{C:\mcp-deploy\matlab}:
  \begin{enumerate}[label=(\alph*)]
    \item Open \url{https://github.com/matlab/simulink-agentic-toolkit}.
    \item Click \textbf{Code} \textrightarrow{} \textbf{Download ZIP}.
    \item Extract it the same way, setting the destination to \path{C:\mcp-deploy\matlab} itself.
    \item \textbf{Rename} the extracted \path{simulink-agentic-toolkit-main} folder to
          \path{simulink-agentic-toolkit}.
  \end{enumerate}
\end{enumerate}

% ============================================================
\section{Raise the Context Limit (Optional)}
% ============================================================

Also identical either way --- there is no one universal setting name, so the same general
guidance applies regardless of configuration format. Do this if the bridge and model in use
allow it; skip it otherwise and come back to it later if tool calls start failing.

\begin{enumerate}
  \item If a context-length or token-limit setting is exposed for whatever is actually running
        the model underneath the bridge, raise it comfortably above what a full tool set plus
        skill text is expected to need.
\end{enumerate}

\begin{notebox}
  A too-small context limit here will \textbf{silently truncate} tool definitions and skill
  content instead of erroring cleanly --- err generous if adjustable, or shrink the tool/skill
  selection if not.
\end{notebox}


% ============================================================
\section{Register the MCP Server}
% ============================================================

\methodA

\begin{enumerate}
  \item Open the bridge's TOML config file --- commonly \path{config.toml} in a hidden
        dot-folder in your home directory.

  \item Add this block:

  \begin{codeblock}
  \begin{Verbatim}
[[mcp_servers]]
name = "matlab-mcp"
transport = "stdio"
command = "C:/mcp-deploy/matlab/matlab-mcp-server/matlab-mcp-server-windows-x64.exe"
args = ["--extension-file=C:/mcp-deploy/matlab/simulink-agentic-toolkit/tools/tools.json"]
  \end{Verbatim}
  \end{codeblock}
\end{enumerate}

\methodB

\begin{enumerate}
  \item Locate the bridge's config file that lists MCP servers --- commonly JSON, with an
        \path{mcpServers} object.

  \item Back up the file.

  \item Add this entry, or copy a starting template from
        \path{matlab-agentic-toolkit/templates/}:

  \begin{codeblock}
  \begin{Verbatim}
{
  "mcpServers": {
    "matlab-mcp": {
      "command":
        "C:\\mcp-deploy\\matlab\\matlab-mcp-server\\matlab-mcp-server-windows-x64.exe",
      "args": [
        "--extension-file=C:\\mcp-deploy\\matlab\\simulink-agentic-toolkit\\tools\\tools.json"
      ]
    }
  }
}
  \end{Verbatim}
  \end{codeblock}
\end{enumerate}


% ============================================================
\section{Place the Skills}
% ============================================================

\methodA

\begin{enumerate}
  \item Find the bridge's skills folder --- check its documentation for the exact location, or
        look for a folder named \path{.agents/skills/} (project-local) or
        \path{~/.agents/skills/} (user-global) in your home directory. Depending on the bridge,
        this folder may not exist yet --- create it if it doesn't.

  \item \textbf{Set \texttt{\$dst} below to that skills folder}, then run this in PowerShell ---
        \path{skills-catalog} nests each skill inside a category folder, so this finds every
        \path{SKILL.md} wherever it sits and copies just that skill's folder straight into
        \texttt{\$dst}, for both toolkits in one pass:

  \begin{codeblocknobreak}
  \begin{Verbatim}
$dst = "<bridge's skills folder>"
$srcs = "C:\mcp-deploy\matlab\matlab-agentic-toolkit\skills-catalog",
        "C:\mcp-deploy\matlab\simulink-agentic-toolkit\skills-catalog"
foreach ($src in $srcs) {
  if (Test-Path $src) {
    Get-ChildItem $src -Directory -Recurse | Where-Object {
      Test-Path (Join-Path $_.FullName "SKILL.md")
    } | ForEach-Object {
      Copy-Item $_.FullName -Destination (Join-Path $dst $_.Name) -Recurse -Force
    }
  }
}
  \end{Verbatim}
  \end{codeblocknobreak}

        The second toolkit path is simply skipped if it was not downloaded.

  \item If the bridge's config file has an \path{enabled_skills} list, add the copied skill
        names to it, then reload or restart.
\end{enumerate}

\methodB

\begin{itemize}
  \item \textbf{If it expects skills at a fixed path:} set \texttt{\$dst} to that path and run
        the same PowerShell command shown above under Method A --- it copies every skill folder
        from both toolkits in one pass and simply skips whichever toolkit was not downloaded.
\end{itemize}


% ============================================================
\section{Restart and Verify}
% ============================================================

\begin{enumerate}
  \item Restart or reload the bridge, whichever it supports, so the configuration change takes
        effect.
  \item Confirm the MATLAB server appears in the bridge's tool list, with its full tool set, not
        a truncated one.
  \item Ask your agent:
\end{enumerate}

\begin{codeblock}
\begin{Verbatim}
What version of MATLAB is running? List the installed toolboxes.
\end{Verbatim}
\end{codeblock}


\vfill
\begin{center}
  {\color{TSgray}\scriptsize\textit{\copyright{} \the\year{} TechSource. All rights reserved.}}
\end{center}

\end{document}
